% !TeX root = ../main.tex

\chapter{Dataverwerking en lokale opslag}
\label{chap:data}

\pending{Dit hoofdstuk inhoudelijk nakijken en afwerken. Controleer de technische details over tijdseenheden, providersemantiek, ronde-identificatie, bestandsformaten en herstel na een herstart.}

\section{Uniforme parser}
De verschillende timingproviders gebruiken verschillende labels en representaties. De parser begint daarom met het
opschonen van witruimte en het canoniseren van headers. Varianten worden gekoppeld aan interne
sleutels voor positie, startnummer, klasse, PIC (Place In Class), team, piloot, aantal ronden, gap, interval,
sectoren, laatste tijd, beste tijd en pitinformatie. Een headercontrole voorkomt dat een
willekeurige tabel op de pagina als timingtabel wordt gebruikt.

Tijdswaarden worden zo vroeg mogelijk naar milliseconden omgezet. De tekst \code{2:05.073} wordt
intern bijvoorbeeld \code{125.073 milliseconden}. Formatteringsfuncties zetten dit pas bij presentatie
opnieuw om. Dit maakt optellen, gemiddeldes nemen, vergelijken en testen eenvoudiger.

Elke genormaliseerde rij bevat ook de gedetecteerde bronprovider. Dat is noodzakelijk omdat een
identiek ogend veld een andere semantiek kan hebben. GetRaceResults levert een bruikbare \emph{DIFF}
tussen opeenvolgende algemene posities. Om de afstand tussen twee klassewagens te bepalen worden
alle tussenliggende DIFF-waarden opgeteld, ook die van wagens uit een andere klasse. RIS Timing
gebruikt \emph{INT}; wanneer die waarde een ronde-eenheid bevat, wordt de rondeachterstand
conservatief omgerekend met een representatieve gemiddelde rondetijd.

\begin{comment}
%%% Herschrijf ofzo? Is dit nodig? TODO
\section{Detectie van voltooide ronden}
De live tabel bevat zowel de laatst voltooide rondetijd als sectoren van de ronde die bezig is. Bij
elke poll wordt de laatste bekende rij per wagen bewaard. Wanneer de combinatie van wagen, ronde en
tijd een nieuwe voltooide ronde aanduidt, wordt een laprecord gemaakt. De sectorwaarden van de
vorige snapshot vormen dan het beste bewijs voor de sectoren van die voltooide ronde.

Een lap identity voorkomt duplicaten wanneer dezelfde live rij gedurende meerdere polls zichtbaar
blijft. Bij een herstart worden identities uit de bestaande JSONL-historiek opnieuw in een set
geladen. Nieuwe data kan daardoor veilig aan dezelfde sessie worden toegevoegd zonder oude ronden
opnieuw te schrijven.
\end{comment}


\section{Vlagstatus en pitfasen}
Voor elke sector wordt de vlagstatus vastgelegd op het moment dat de sector beschikbaar wordt. Dit
detail is nodig voor het scenario waarin sector 1 groen werd gereden en tijdens sector 2 FCY
begint. De volledige ronde is dan niet representatief voor een rondetijdgemiddelde, maar sector 1
kan nog geldig zijn voor het gemiddelde en de beste waarde van sector 1.

De analyselaag annoteert ook pitfasen. Veranderingen in pitteller en pitstatus helpen inlaps en
outlaps identificeren. Deze ronden blijven deel van de historische werkelijkheid, maar worden niet
gebruikt voor representatief tempo. Dit voorkomt dat een trage pitronde het gemiddelde van een
piloot kunstmatig verslechtert.

\section{Bestandsformaat}
Er werd uiteindelijk niet voor SQLite gekozen. Voor de schaal van één race en de behoefte aan
inspecteerbaarheid bleken leesbare bestanden eenvoudiger en voldoende. De belangrijkste bestanden
zijn:

\begin{itemize}
  \item \code{lap_history.jsonl}: append-only historiek met 1 JSON-object per opgeslagen ronde
  \item \code{lap_history.csv}: dezelfde rondehistoriek in tabelvorm voor manuele inspectie of spreadsheetanalyse
  \item \code{latest_live_rows.json} en \code{latest_live_rows.csv}: laatste volledige timingsnapshot
  \item \code{session_metadata.json}: sessiecontext, gevolgde wagen(s), timingprovider en instellingen
  \item \code{parser_debug.json}: gedetecteerde headers, voorbeelden en parserwaarschuwingen
  \item \code{analytics_summary.json}: laatst berekende statistieken en aggregaten voor dashboard en grafieken
  \item \code{lap_prediction.json} en \code{lap_prediction_car-<nummer>.json}: actuele rondevoorspelling per gevolgde wagen
  \item \code{pitstop_plan.json} en \code{pitstop_plan_car-<nummer>.json}: actuele pitstopplanning per gevolgde wagen
  \item \code{gap_state.json} en \code{gap_history.jsonl}: laatst bevestigde gapinformatie en gap-historiek
  \item \code{stint_state.json}: actuele stintinformatie per wagen en piloot
  \item \code{track_condition_state.json} en \code{track_condition_events.jsonl}: huidige baanconditie en wijzigingen doorheen de sessie
\end{itemize}

JSONL is geschikt voor rondes omdat een nieuwe ronde kan worden toegevoegd zonder het volledige
bestand te herschrijven. Tegelijk blijft elke regel leesbaar en herstelbaar. Samenvattingsbestanden
worden wel overschreven: ze vertegenwoordigen de actuele afgeleide toestand en kunnen indien nodig
opnieuw uit de lap history worden berekend.

\begin{comment}
\section{Incrementele statistiek}
De applicatie bewaart naast ruwe historiek ook de laatste analysesamenvatting. Dit maakt inspectie
eenvoudig, maar de lap history blijft de gezaghebbende bron. Gemiddelden worden bij een update
opnieuw opgebouwd uit de relevante gefilterde ronden. Voor de verwachte omvang van maximaal enkele
duizenden ronden is dit computationeel klein en betrouwbaarder dan complexe incrementele correcties.
Wanneer een ronde later als pit- of neutralisatieronde wordt herkend, zou een uitsluitend
incrementeel gemiddelde immers ook moeten worden teruggedraaid.
\end{comment}


\section{Sessiemappen en herstel}
De gebruiker selecteert voor elke sessie een nieuwe map. Daardoor kan dezelfde map na een crash of
foutieve sluiting opnieuw worden geselecteerd. Bij het starten laadt de applicatie bestaande lap
history en pitstatus, bouwt de duplicate guard op en vervolgt de collectie.
